\documentclass[12pt]{article}
\usepackage{mathpazo}
\usepackage[T1]{fontenc}
\usepackage[utf8]{inputenc}
\usepackage{geometry}
\geometry{verbose,tmargin=1cm,bmargin=2.7cm,lmargin=2.7cm,rmargin=2cm}
\pagestyle{empty}
\usepackage{setspace}
\usepackage{amsmath,amssymb}
\usepackage{booktabs}
\usepackage{tabularx}
\usepackage{enumitem}
\setlist[enumerate]{leftmargin=1cm}
\setlist[itemize]{leftmargin=0.5cm}
\onehalfspacing
\setlength{\parskip}{\smallskipamount}
\setlength{\parindent}{0pt}
\setlength{\emergencystretch}{3em}
\date{}
\author{}
\newcommand{\hint}[1]{\par\smallskip{\small\textit{Hint:} #1}\par}
\newcommand{\solution}[1]{}

\begin{document}

\title{\textsc{Assignment II: Rate Cuts versus Stimulus Payments}}
\maketitle
\vspace{-3mm}\thispagestyle{empty}

\textbf{Objective:} This project teaches you when uniform transfers can give the same output path as a rate cut, and why the two policies help different people. It follows Wolf (2025), \emph{Interest Rate Cuts versus Stimulus Payments: An Equivalence Result}, Journal of Political Economy. You do not need to read the paper: all you need is in this document.
\begin{itemize}
\item \textbf{Problem:} The problem consists of
\begin{enumerate}
\item A number of questions (page 2)
\item A model (page 5 onward, incl. calibration and numerical implementation)
\item Some code files you can start from
\end{enumerate}
\item \textbf{Code:} The problem is designed to be solved with the \texttt{GEModelTools} package.
\item \textbf{Figures and tables:} All plotting and printing code is given in the notebook. Use the provided plotting and printing functions for all figures and tables.
\item \textbf{Structure:} Your project should consist of
\begin{enumerate}
\item A single self-contained pdf-file with all results
\item A single Jupyter notebook showing how the results are produced
\item Well-documented \texttt{.py} files
\end{enumerate}
\item \textbf{Hand-in:} Upload a single zip-file on Absalon (and nothing else)
\item \textbf{Deadline:} 24 November 2026
\item \textbf{Exam:} This project will be a part of your exam portfolio.
You can incorporate feedback before handing in the final version.
\end{itemize}

\newpage

\section*{Questions}

The two experiments are defined in Section 1. The \emph{rate cut} is a surprise cut of 25 basis points in the real rate, with persistence $\rho_r=0.4$. The \emph{transfer shock} is a path $d\boldsymbol{\chi}$ of uniform transfers, announced at $t=0$, with the rate constant.

\begin{enumerate}
\item \textbf{The Intertemporal Keynesian Cross (IKC).} These questions are pen and paper.
\begin{enumerate}
\item Combine the goods market, the consumption function \eqref{eq:cons}, $d\boldsymbol{Z}=d\boldsymbol{Y}-d\boldsymbol{T}$ and the tax rule \eqref{eq:taxes} to show that, for a rate shock, the output response is given by
\begin{equation}
(\boldsymbol{I}-\boldsymbol{M})d\boldsymbol{Y}=\left(\boldsymbol{M}^{r}-B_{ss}\boldsymbol{M}\boldsymbol{L}\right)d\boldsymbol{r},\label{eq:ikc_monetary}
\end{equation}
and for a transfer shock, the output response is given by
\begin{equation}
(\boldsymbol{I}-\boldsymbol{M})d\boldsymbol{Y}=\boldsymbol{M}^{\chi}d\boldsymbol{\chi}.\label{eq:ikc_transfer}
\end{equation}
\solution{1a}
\item Explain each equation in words. What is the effect due to $\boldsymbol{M}$ on the left-hand side? Why does $B_{ss}\boldsymbol{M}\boldsymbol{L}$ appear? What is the difference between $\boldsymbol{M}$ and $\boldsymbol{M}^{\chi}$?
\solution{1b}
\item Assume that each IKC has a unique solution that returns to steady state, and that $\boldsymbol{M}^{\chi}$ can be inverted. Show that the rate cut and the transfer shock give the same $d\boldsymbol{Y}$ if
\begin{equation}
d\boldsymbol{\chi}=(\boldsymbol{M}^{\chi})^{-1}\left(\boldsymbol{M}^{r}-B_{ss}\boldsymbol{M}\boldsymbol{L}\right)d\boldsymbol{r}.\label{eq:equiv}
\end{equation}
Use the W-NKPC \eqref{eq:nkwc} to show that inflation is then the same as well.
\solution{1c}
\item Note that $\boldsymbol{q}'\boldsymbol{M}^{r}d\boldsymbol{r}=B_{ss}\boldsymbol{q}'\boldsymbol{L}d\boldsymbol{r}$ and $\boldsymbol{q}'\boldsymbol{M}=\boldsymbol{q}'\boldsymbol{M}^{\chi}=\boldsymbol{q}'$, as shown in Section 1. Show that the transfers computed with \eqref{eq:equiv}, which reproduce the output path of a rate cut $d\boldsymbol{r}$, have zero net present value: $\boldsymbol{q}'d\boldsymbol{\chi}=0$.
\solution{1d}
\end{enumerate}

\item \textbf{RANK.} We now ask whether \eqref{eq:equiv} can be used in RANK, defined at the end of Section 1. In RANK, a transfer is income like any other, so $\boldsymbol{M}^{\chi}=\boldsymbol{M}=\boldsymbol{M}^{RA}$, as derived in the lecture on fiscal policy. Since $\beta(1+r_{ss})=1$ in RANK, $q_{s}=\beta^{s}$.
\begin{enumerate}
\item \emph{The transfer shock.} Recall from the lecture that
$$M^{RA}_{t,s}=(1-\beta)\beta^{s}\qquad\text{for all }t,\qquad\text{so}\qquad\boldsymbol{M}^{RA}=(1-\beta)\boldsymbol{1}\boldsymbol{q}'.$$
Show that $\boldsymbol{M}^{\chi}d\boldsymbol{\chi}=\boldsymbol{0}$ for any transfer shock with zero present value, and use \eqref{eq:ikc_transfer} to find the output response. Assume that output returns to steady state.
\solution{2a}
\item \emph{The rate cut.} Use the Euler equation \eqref{eq:euler} to show that output after the rate cut is
$$dY_t=-\beta\frac{C_{ss}}{\sigma}\sum_{k=t}^{\infty}dr_k.$$
How much does output rise on impact?
\hint{Use $C_{ss}=Y_{ss}-G$ and the geometric sum of $dr_k$.}
\solution{2b}
\item \emph{Conclude.} Can a transfer shock with zero net present value reproduce the output response to a rate cut in RANK? Why?
\solution{2c}
\end{enumerate}

\item \textbf{HANK.} The model code in \texttt{blocks.py} is given. Read it and check that it matches Section 1.
\begin{enumerate}
\item \emph{The Jacobians.} Compute the Jacobians $\boldsymbol{M}$, $\boldsymbol{M}^{\chi}$ and $\boldsymbol{M}^{ra}$ (see Section 3), and plot their columns $s=0,5,10,20$. How do $\boldsymbol{M}^{\chi}$ and $\boldsymbol{M}$ differ from RANK?
\solution{3a}
\item \emph{The rate cut.} Compute the output response to the rate cut with \texttt{find\_IRFs}, and check that the IKC \eqref{eq:ikc_monetary} gives the same response.
\hint{Check: output rises by 0.173\% on impact.}
\solution{3b}
\item \emph{The transfer shock.} Compute the transfer path $d\boldsymbol{\chi}$ from \eqref{eq:equiv}. Then, compute the output response to this shock, and show that output and inflation coincide with the rate cut. Why does this work in HANK and not in RANK?
\solution{3c}
\item \emph{The size of the transfers.} How large is the transfer shock, and how large is the rise in debt? In which year does $d\chi_{t}$ turn negative? Check that the transfer shock has zero present value.
\solution{3d}
\end{enumerate}

\item \textbf{Welfare.} The two shocks give the same output. We now ask who benefits. To keep things simple, we measure welfare with the utility of consumption only, and ignore the disutility of work. The household problem returns $u_{t}=\frac{c_{t}^{1-\sigma}}{1-\sigma}$, and its aggregate $U_{t}=\int u_{t}d\mathcal{D}_{t}$ is \texttt{U\_hh}. The welfare of a group of households with mass one, who start at $t=0$ with the distribution $\mathcal{D}_{0}$, is
$$W=\sum_{t=0}^{T-1}\beta^{t}U_{t},$$
where $\mathcal{D}_{t}$ follows the group forward in time. Let $W_{ss}$ be its welfare when all inputs stay at steady state and $W_{0}$ its welfare after a shock. The consumption equivalent $\lambda$ is the permanent rise in consumption, as a fraction (the notebook reports it in percent), that gives the group the same welfare in steady state as after the shock:
$$(1+\lambda)^{1-\sigma}W_{ss}=W_{0}.$$
Complete \texttt{welfare} in the notebook, and use the provided code to print the welfare change $\lambda$ from the rate cut and the transfer shock by wealth quintile, from Q1 (the poorest) to Q5 (the richest). Who benefits most from each shock, who benefits least, and why? \texttt{model.decompose\_hh\_path} solves and simulates the households from a given initial distribution \texttt{Dbeg}, along the paths of inputs given in \texttt{custom\_paths}. The notebook builds \texttt{custom\_paths} for you: the paths are in levels, $X_{ss}+dX_{t}$, for the three household inputs, \texttt{Z}, \texttt{ra} and \texttt{chi}. List the same three inputs in \texttt{use\_inputs}. \texttt{path.U\_hh} has shape $(T,1)$.
\hint{Check: after the rate cut, $\lambda$ is 0.053\% in Q1 and 0.014\% in Q5. Think about who receives what, and who pays later.}
\solution{4}

\item \textbf{Conclusion.} Suppose that the central bank cannot cut the rate because it is at a lower bound. Based on your results, can transfers replace the rate cut? Outside a lower bound, which of the two would you use, and why? What does the model leave out that could change your answer?
\solution{5}

\item \textbf{Capital gains.} In the current model, we have abstracted from capital gains by assuming no price markup, $\mu=1$. How do you think capital gains would affect the distribution of welfare following a rate cut? Who would benefit from capital gains? No code is needed.
\solution{6}

\end{enumerate}


\newpage

\section{Model}

The model is the HANK model of the lectures on fiscal and monetary policy, with a real rate rule, real bonds and a uniform transfer. Time is discrete and annual, as in the lectures. A bold letter is a path, and $d\boldsymbol{x}=(dx_0,dx_1,\ldots)'$ is the deviation of the path from steady state. $\boldsymbol{L}$ is the lag matrix, so $(\boldsymbol{L}d\boldsymbol{x})_t=dx_{t-1}$, and $\boldsymbol{q}$ is the vector of discount factors, $q_s=(1+r_{ss})^{-s}$. All responses in Questions 1 to 3 are linear.

\paragraph*{Households.}
Households have idiosyncratic productivity $z_{t}$ and save in real government bonds:
\begin{align*}
v_{t}(z_{t},a_{t-1}) & =\max_{c_{t}}\frac{c_{t}^{1-\sigma}}{1-\sigma}-\varphi\frac{\ell_{t}^{1+\nu}}{1+\nu}+\beta\mathbb{E}_{t}\left[v_{t+1}(z_{t+1},a_{t})\right]\\
\text{s.t. }a_{t}+c_{t} & =(1+r_{t-1})a_{t-1}+(1-\tau_{t})w_{t}\ell_{t}z_{t}+\chi_{t}\\
\log z_{t+1} & =\rho_{z}\log z_{t}+\psi_{t+1},\quad\psi_{t}\sim\mathcal{N}(\mu_{\psi},\sigma_{\psi}^{2}),\quad\mu_{\psi}\text{ such that }\mathbb{E}[z_{t}]=1\\
a_{t} & \geq0
\end{align*}
$r_{t}$ is the real rate from $t$ to $t+1$: it is known at $t$ and paid at $t+1$. Hours $\ell_{t}=L_{t}$ are set by the union and are the same for all households, so they are not a choice. $\tau_{t}$ is the labor tax rate and $\chi_{t}$ a transfer, the same for all households. Aggregate consumption and assets are $C_{t}^{hh}=\int c_{t}d\mathcal{D}_{t}$ and $A_{t}^{hh}=\int a_{t}d\mathcal{D}_{t}$.

The consumption function is $\boldsymbol{C}^{hh}=C^{hh}(\boldsymbol{Z},\boldsymbol{r},\boldsymbol{\chi})$, with post-tax labor income $Z_{t}=(1-\tau_{t})w_{t}L_{t}$. To first order,
\begin{equation}
d\boldsymbol{C}^{hh}=\boldsymbol{M}d\boldsymbol{Z}+\boldsymbol{M}^{r}d\boldsymbol{r}+\boldsymbol{M}^{\chi}d\boldsymbol{\chi}.\label{eq:cons}
\end{equation}
$\boldsymbol{M}$ is the matrix of intertemporal MPCs (iMPCs) of the lectures, $M_{t,s}=\partial C_{t}^{hh}/\partial Z_{s}$, and $\boldsymbol{M}^{r}$ is the Jacobian with respect to $\boldsymbol{r}$. $\boldsymbol{M}^{\chi}$ is new: $M^{\chi}_{t,s}=\partial C_{t}^{hh}/\partial\chi_{s}$ is the change in consumption at date $t$ after a change in the transfer at date $s$. It differs from $\boldsymbol{M}$ because $\chi_{t}$ is the same for all households, while $Z_{t}$ is scaled by $z_{t}$. Its first entry $M^{\chi}_{0,0}$ is the MPC on impact: the share of a transfer paid at $t=0$ that is spent at $t=0$.

\paragraph*{Present values.}
The household budget pins down the present value of the consumption response. Summed over households, the budget is $C_{t}^{hh}+A_{t}^{hh}=(1+r_{t-1})A_{t-1}^{hh}+Z_{t}+\chi_{t}$. Change only $\boldsymbol{r}$ and linearize:
$$dC_{t}^{hh}+dA_{t}^{hh}=(1+r_{ss})dA_{t-1}^{hh}+A_{ss}\,dr_{t-1}.$$
Multiply by $q_{t}$ and sum over $t$. Since $q_{t}(1+r_{ss})=q_{t-1}$, the asset terms cancel, because $dA_{-1}^{hh}=0$ and assets return to steady state. With $A_{ss}=B_{ss}$, this gives
\begin{equation}
\boldsymbol{q}'\boldsymbol{M}^{r}d\boldsymbol{r}=B_{ss}\boldsymbol{q}'\boldsymbol{L}d\boldsymbol{r}:\label{eq:pv_r}
\end{equation}
in present value, households spend exactly the change in their interest income. The same steps with a change in $Z_{s}$ or $\chi_{s}$ give $\boldsymbol{q}'\boldsymbol{M}=\boldsymbol{q}'\boldsymbol{M}^{\chi}=\boldsymbol{q}'$: a unit of income is eventually spent in full.

\paragraph*{Firms and unions.}
Firms produce with labor only, $Y_{t}=L_{t}$, with no price markup, $\mu=1$, so $w_{t}=1$ and $\pi_{t}=\pi^{w}_{t}$. Unions set wages as in the lectures, with $h(\ell)=\varphi\frac{\ell^{1+\nu}}{1+\nu}$ and $u(c)=\frac{c^{1-\sigma}}{1-\sigma}$. This gives the wage Phillips curve (W-NKPC)
\begin{equation}
\pi^{w}_{t}=\kappa\left(h'\left(L_{t}\right)-\frac{1-\tau_{ss}}{\mu^{w}}u'\left(C_{t}^{hh}\right)\right)+\beta\pi^{w}_{t+1},\label{eq:nkwc}
\end{equation}
and households work the hours that firms demand. The tax wedge is kept at its steady-state value, so the W-NKPC depends only on $L_{t}$ and $C_{t}^{hh}$. With $\tau_{t}$ instead, it would also depend on the tax rate, which moves after a rate cut but not after a transfer shock.

\paragraph*{Central bank.}
The central bank sets the real rate, $r_{t}=r_{ss}+\epsilon^{r}_{t}$. The \emph{rate cut} is $\epsilon^{r}_{0}=-0.0025$ with $\epsilon^{r}_{t}=\rho_{r}\epsilon^{r}_{t-1}$. Since bonds pay the rate set in the previous period, the return at $t=0$ does not move.

\paragraph*{Government.}
Government spending $G$ is constant. The government has two instruments. Labor taxes $T_{t}=\tau_{t}Y_{t}$ cover spending and the interest on steady-state debt,
\begin{equation}
T_{t}=G+r_{t-1}B_{ss},\quad\text{so}\quad d\boldsymbol{T}=B_{ss}\boldsymbol{L}d\boldsymbol{r}.\label{eq:taxes}
\end{equation}
As in the lecture on monetary policy, a rate cut lowers the interest bill, and taxes fall with it. Transfers $\chi_{t}$ are a path chosen at $t=0$, paid for with debt:
\begin{equation}
B_{t}=(1+r_{t-1})B_{t-1}+G+\chi_{t}-T_{t},\quad\text{so to first order}\quad dB_{t}=(1+r_{ss})dB_{t-1}+d\chi_{t}.\label{eq:gov}
\end{equation}
After the rate cut, debt is constant. Debt returns to steady state if the transfers have zero present value, $\boldsymbol{q}'d\boldsymbol{\chi}=0$. Post-tax labor income is $Z_{t}=Y_{t}-T_{t}$.

\paragraph*{Market clearing.}
The asset market clears, $B_{t}=A^{hh}_{t}$, and the goods market clears, $Y_{t}=C^{hh}_{t}+G$.

\paragraph*{RANK.}
In RANK, a representative household with the same preferences, $\beta(1+r_{ss})=1$ and no borrowing constraint replaces the households above. It solves
\begin{align*}
\max_{C_{t},A_{t}} & \sum_{t=0}^{\infty}\beta^{t}\left[u\left(C_{t}\right)-h\left(L_{t}\right)\right]\\
\text{s.t. }C_{t}+A_{t} & =\left(1+r_{t-1}\right)A_{t-1}+Z_{t}+\chi_{t},
\end{align*}
which gives the Euler equation
\begin{equation}
u'\left(C_{t}\right)=\beta\left(1+r_{t}\right)u'\left(C_{t+1}\right).\label{eq:euler}
\end{equation}
Everything else is unchanged.

\section{Calibration}

\begin{table}[h]
\centering
\caption{Parameters}
\small
\begin{tabularx}{\textwidth}{lXl}
\toprule
 & Parameter & Value \\
\midrule
$\sigma$ & inverse elasticity of intertemporal substitution & 2.0 \\
$\nu$ & inverse Frisch elasticity & 1.0 \\
$\rho_z$, $\sigma_\psi$ & income process & $0.95^{4}$, 0.186 \\
$\mu^{w}$ & wage markup & 1.2 \\
$\kappa$ & slope of the W-NKPC & 0.16 \\
$G$ & government spending & 0.20 \\
$B_{ss}$ & government bonds, annual output & 0.25 \\
$r_{ss}$ & steady-state real rate & 0.02 \\
$\rho_r$ & persistence of the rate cut & 0.4 \\
$\beta$ & discount factor & asset market clearing \\
$\varphi$ & disutility of work & $L=1$ in steady state \\
$\tau_{ss}$ & labor tax rate in steady state & government budget \\
\bottomrule
\end{tabularx}
\end{table}

There are no transfers in steady state, $\chi_{ss}=0$, and $Y_{ss}=1$.

\section{Numerical implementation}

\begin{itemize}
\item The code is the HANK code of the lectures. The two shocks are \texttt{eps\_r} and \texttt{chi}.
\item As in the lecture code, the household input is the return paid at $t$, \texttt{ra}, so \texttt{ra[t]} is $r_{t-1}$, and \texttt{r[t]} is $r_{t}$.
\item The Jacobians are in \texttt{model.jac\_hh}: \texttt{model.jac\_hh[('C\_hh','Z')]} is $\boldsymbol{M}$ and \texttt{model.jac\_hh[('C\_hh','chi')]} is $\boldsymbol{M}^{\chi}$. \texttt{model.jac\_hh[('C\_hh','ra')]} is $\boldsymbol{M}^{ra}$, the Jacobian with respect to the return paid at $t$, $d\boldsymbol{r}^{a}=\boldsymbol{L}d\boldsymbol{r}$. In the code, work with \texttt{dra} $=d\boldsymbol{r}^{a}$ instead of $d\boldsymbol{r}$: $\boldsymbol{M}^{r}d\boldsymbol{r}=\boldsymbol{M}^{ra}d\boldsymbol{r}^{a}$ and $\boldsymbol{L}d\boldsymbol{r}=d\boldsymbol{r}^{a}$, so no lag matrix is needed.
\item Linear responses come from \texttt{find\_IRFs}. The household problem also returns the utility of consumption \texttt{u}, aggregated to \texttt{U\_hh}. Question 4 feeds the responses into the household problem to get the welfare of each wealth quintile.
\item The horizon is $T=500$. Productivity is discretized with the Rouwenhorst method with 7 states. The asset grid has 500 points between 0 and 50.
\end{itemize}


\end{document}
